Fix a sufficiently large index for which \(\mathcal M_{n,d,L}\ne\varnothing\), and select the graph primitives \(W\), \((\lambda_k,\psi_k)_{k=1}^r\), and graph sampling law from one admissible member.  We retain those primitives and change only the response law.  Constants below may depend on the fixed class constants and \(\alpha\), but never on \(n\) or \(d\).  Since \(C_{\log}\log n\leq d\) with \(C_{\log}>0\), we have \(d\to\infty\); hence \(\mathcal H_{n,d}\ne\varnothing\) for every sufficiently large index, so all criteria and envelopes used below are defined under their stated nonempty-domain conventions.

Define
\[
 w(u)=\int_0^1W(u,v)\,dv,
 \qquad p_n(u)=\rho_nw(u).
\tag{1}
\]
By \({\rm ass:kernel\text{-}bound}\), \({\rm ass:degree\text{-}regularity}\), and \({\rm ass:edge\text{-}probability}\), outside a Lebesgue-null set,
\[
 \kappa^{-1}\leq w(u)\leq\kappa,
 \qquad 0<p_n(u)\leq1.
\tag{2}
\]
The uniform marginal in \({\rm ass:unit\text{-}sampling}\) avoids the exceptional set almost surely.  Conditional on \(U_i=u\), integrating \({\rm ass:graphon\text{-}draw}\) over the independent uniform variables \(U_j\) gives independent edges with common success probability \(p_n(u)\).  Therefore
\[
 N_i\mid U_i=u\sim{\rm Binomial}(n-1,p_n(u)),
 \qquad m_n(u):=\mathbb E(N_i\mid U_i=u)=(n-1)p_n(u),
\tag{3}
\]
and, because \(\rho_n=d/n\),
\[
 \frac{(n-1)d}{n\kappa}\leq m_n(u)
 \leq\frac{(n-1)d\kappa}{n}.
\tag{4}
\]
Thus \(d/(2\kappa)\leq m_n(u)\leq\kappa d\) for \(n\geq2\).

We record the degree consequences used below.  If \(N\sim{\rm Binomial}(n-1,p)\), direct expansion of its mass function and the identity \({n-1\choose k}/(k+1)={n\choose k+1}/n\) give
\[
 \mathbb E\frac1{N+1}
 =\frac{1-(1-p)^n}{np}\leq\frac1{np}.
\tag{5}
\]
Since \(N^{-1}\mathbf1\{N>0\}\leq2(N+1)^{-1}\), (2)--(5) imply
\[
 \mathbb E\{N_i^{-1}\mathbf1(N_i>0)\mid U_i=u\}
 \leq\frac{2\kappa}{d}.
\tag{6}
\]
Also,
\[
 \Pr(N_i=0\mid U_i=u)=(1-p_n(u))^{n-1}
 \leq e^{-m_n(u)}\leq e^{-d/(2\kappa)}.
\tag{7}
\]
Conditional on \(A\) and \(N_i=N>0\), \({\rm ass:bernoulli\text{-}design}\) gives \(M_i\sim{\rm Binomial}(N,1/2)\), so
\[
 \mathbb E_Z(V_i^2\mid A,N_i=N)=\frac1{4N}.
\]
Integrating and using (6) yields
\[
 \mathbb E\sum_{i=1}^nV_i^2\mathbf1\{N_i>0\}
 \leq\frac{n\kappa}{2d}.
\tag{8}
\]

For \(M\sim{\rm Binomial}(N,1/2)\), its largest atom is at most \(C_0N^{-1/2}\).  Indeed, for \(a_q={2q\choose q}4^{-q}\),
\[
 \frac{a_{q+1}}{a_q}=\frac{2q+1}{2q+2},
 \qquad
 \left(\frac{2q+1}{2q+2}\right)^2\leq\frac{q+1}{q+2};
\]
induction gives \(a_q\leq(q+1)^{-1/2}\), and the odd case is bounded by an adjacent even central atom.  The interval \(|M/N-1/2|\leq h\) contains at most \(2Nh+2\) lattice points, whence
\[
 \Pr\{|M/N-1/2|\leq h\mid N\}
 \leq C_1(\sqrt N\,h+N^{-1/2}).
\tag{9}
\]
For a binomial variable \(N\) of mean \(m\), exponential Markov's inequality and the binomial moment-generating function give, for every \(t>0\),
\[
 \Pr(N<m/2)
 \leq e^{tm/2}\mathbb E e^{-tN}
 \leq\exp\{tm/2-m(1-e^{-t})\}.
\]
Taking \(t=\log2\) gives
\[
 \Pr(N<m/2)\leq e^{-c_0m},
 \qquad c_0=(1-\log2)/2>0.
\tag{10}
\]
Splitting at \(N_i=m_n(U_i)/2\), and using (4), (10), and \(d\to\infty\), gives
\[
 \mathbb E\{N_i^{-1/2}\mathbf1(N_i>0)\}
 \leq C_2d^{-1/2}.
\tag{11}
\]
Jensen's inequality and (4) give \(\mathbb E\sqrt{N_i}\leq\sqrt{\kappa d}\).  Combining this with (9), (11), and \(h\geq d^{-1}\) yields, uniformly over \(h\in\mathcal H_{n,d}\),
\[
 \Pr(|V_i|\leq h,N_i>0)\leq C_3\sqrt d\,h.
\tag{12}
\]

Choose \(K>e\kappa\) so large that
\[
 a_K:=K\log\{K/(e\kappa)\}>2/C_{\log},
 \qquad D_n=\lceil Kd\rceil.
\tag{13}
\]
Optimizing exponential Markov's inequality for the binomial law in (3) gives
\[
 \Pr(N_i>D_n\mid U_i=u)
 \leq\left\{\frac{em_n(u)}{D_n}\right\}^{D_n}
 \leq e^{-a_Kd}.
\]
No independence among degrees is needed for the union bound, so \({\rm ass:log\text{-}degree}\) implies
\[
 \Pr\left(\max_iN_i>D_n\right)
 \leq ne^{-a_Kd}
 \leq n^{1-a_KC_{\log}}=o(1).
\tag{14}
\]

We now construct the response experiments.  Set \(s_0=B/4>0\), and let \((\xi_i)\), independently of \((U,A,Z)\), be independent with density
\[
 p_{s_0}(t)=s_0^{-1}\cos^2\!\left(\frac{\pi t}{2s_0}\right)
 \mathbf1\{|t|\leq s_0\}.
\tag{15}
\]
The substitution \(y=\pi t/(2s_0)\) shows that (15) integrates to one.  For deterministic \(g\in C^3[0,1]\), define
\[
 f_i^\pm(a,x)=\xi_i\pm g(x).
\tag{16}
\]
Whenever
\[
 \|g\|_\infty\leq3B/4,
 \qquad \max_{1\leq j\leq2}\|g^{(j)}\|_\infty\leq B,
 \qquad \|g^{(3)}\|_\infty\leq L,
\tag{17}
\]
the response envelope and third-derivative assumption hold.  The pairs \((U_i,f_i^\pm)\) remain independent and identically distributed with uniform \(U_i\)-marginal, and the retained graph satisfies all other member properties.  Thus each sign in (16) defines a law in \(\mathcal M_{n,d,L}\).  By \({\rm def:target}\),
\[
 \theta_n^\pm=\pm g'(1/2),
 \qquad \Delta_g:=|\theta_n^+-\theta_n^-|=2|g'(1/2)|.
\tag{18}
\]
Let \(Q_g^\pm\) denote the resulting joint laws of \(D=(A,Z,(Y_i^{\mathrm{obs}})_{i\in\mathcal V_n})\).  For dominated laws with densities \(p,q\), define
\[
 H^2(p,q)=\int(\sqrt p-\sqrt q)^2,
 \qquad \operatorname{TV}(p,q)=\frac12\int|p-q|.
\]
Let \(q_{s_0}=\sqrt{p_{s_0}}\), extended by zero.  It is absolutely continuous because the cosine vanishes at both endpoints, and direct integration gives
\[
 \|q_{s_0}'\|_2^2=\frac{\pi^2}{4s_0^2}.
\tag{19}
\]
For any translations \(\eta_1,\eta_2\), the fundamental theorem of calculus followed by Minkowski's integral inequality gives
\[
 H^2\{p_{s_0}(\mathord\cdot-\eta_1),p_{s_0}(\mathord\cdot-\eta_2)\}
 \leq\frac{\pi^2(\eta_1-\eta_2)^2}{4s_0^2}.
\tag{20}
\]
Put \(X_i^\circ=1/2+V_i\) when \(N_i>0\), and \(X_i^\circ=0\) when \(N_i=0\).  Conditional on \((A,Z)\), the two outcome densities for unit \(i\) are translations whose locations differ by \(2g(X_i^\circ)\).  Thus (20) gives \(H_i^2\leq\pi^2g(X_i^\circ)^2/s_0^2\).  If \(a_i=\int\sqrt{p_{i,+}p_{i,-}}=1-H_i^2/2\), conditional independence and Tonelli's theorem give
\[
 H^2\!\left(\bigotimes_i p_{i,+},\bigotimes_i p_{i,-}\right)
 =2\left(1-\prod_i a_i\right)
 \leq\sum_iH_i^2,
\]
where \(1-\prod_i(1-x_i)\leq\sum_i x_i\) for \(x_i\in[0,1]\).  Since \((A,Z)\) is observed and has the same marginal law under both signs, integration over that common law yields
\[
 H^2(Q_g^+,Q_g^-)
 \leq\frac{\pi^2}{s_0^2}\mathbb E\sum_{i=1}^ng(X_i^\circ)^2.
\tag{21}
\]
Cauchy--Schwarz applied to \(|p-q|=|\sqrt p-\sqrt q|(\sqrt p+\sqrt q)\) gives
\[
 \operatorname{TV}(Q_g^+,Q_g^-)\leq H(Q_g^+,Q_g^-).
\tag{22}
\]

For the graph-floor pair, take
\[
 g_G(x)=c_GB\sqrt{d/n}\,(x-1/2).
\tag{23}
\]
By \({\rm ass:sparse\text{-}limit}\), condition (17) holds for all sufficiently large \(n\) when \(0<c_G\leq1\) is fixed.  Equations (7), (8), and (23) yield
\[
\begin{aligned}
 \mathbb E\sum_i g_G(X_i^\circ)^2
 &=c_G^2B^2\frac dn
 \left\{\mathbb E\sum_iV_i^2\mathbf1(N_i>0)
 +\frac n4\Pr(N_i=0)\right\}\\
 &\leq C_4c_G^2B^2.
\end{aligned}
\tag{24}
\]
Let \(\eta=(1-2\alpha)/2>0\).  Since \(s_0=B/4\), choose a fixed \(c_G>0\) small enough that (21)--(24) give
\[
 \operatorname{TV}(Q_G^+,Q_G^-)\leq\eta,
 \qquad \Delta_G=2c_GB\sqrt{d/n}.
\tag{25}
\]

For the local nonlinear pair suppose \(L>0\), and define
\[
 \phi(t)=t(1-4t^2)^4\mathbf1\{|t|<1/2\}.
\tag{26}
\]
Its fourth-order endpoint zeros make its zero extension a \(C^3(\mathbb R)\) function, and \(\phi'(0)=1\).  Write \(C_j=\|\phi^{(j)}\|_\infty\), and for \(h\in\mathcal H_{n,d}\) set
\[
 g_h(x)=c_ILh^3\phi\!\left(\frac{x-1/2}{h}\right).
\tag{27}
\]
Then \(\|g_h^{(j)}\|_\infty=c_ILh^{3-j}C_j\) for \(0\leq j\leq3\).  Choose \(c_I>0\) so that \(c_IC_3\leq1\), and, if necessary, decrease it so that (17) holds for every sufficiently large \(n\); this is possible because \(h\leq d^{-1/2}\to0\), while \(B,L\) are fixed.  The bump vanishes at the isolate exposure and is supported within \(|V_i|\leq h\).  Hence (12), (21), and (27) give
\[
 H^2(Q_h^+,Q_h^-)
 \leq C_5c_I^2n\sqrt d\,L^2h^7.
\tag{28}
\]
Set
\[
 h_0=\{L^{-2}/(n\sqrt d)\}^{1/7}.
\tag{29}
\]
If \(d^{-1}<h_0<d^{-1/2}\), substituting (29) into (28) and decreasing \(c_I\) if necessary gives
\[
 \operatorname{TV}(Q_{h_0}^+,Q_{h_0}^-)\leq\eta,
 \qquad \Delta_I=2c_ILh_0^2
 =2c_IL^{3/7}(n^2d)^{-1/7}.
\tag{30}
\]

For the lattice pair, continue to suppose \(L>0\), put \(w_n=(4D_n)^{-1}\), and set
\[
 g_D(x)=c_DLw_n^3\phi\!\left(\frac{x-1/2}{w_n}\right),
 \qquad c_DC_3\leq1.
\tag{31}
\]
Because \(D_n\to\infty\), choosing fixed \(c_D>0\) sufficiently small and using the derivative scaling above verifies (17) for all sufficiently large \(n\).  If \(1\leq N_i\leq D_n\), every noncentral exposure point obeys
\[
 |V_i|\geq
 \begin{cases}
 N_i^{-1},&N_i\text{ even},\\
 (2N_i)^{-1},&N_i\text{ odd},
 \end{cases}
 \geq(2D_n)^{-1}>w_n.
\tag{32}
\]
At a central point \(g_D(1/2)=0\) by oddness, and at an isolate \(g_D(0)=0\).  Thus on \(F_n=\{\max_iN_i\leq D_n\}\), the complete observed outcomes under both signs in (16) equal \(\xi_i\).  The restrictions of the two data laws to \(F_n\) agree, and \(F_n\) has the same probability under both.  Therefore (14) gives
\[
 \operatorname{TV}(Q_D^+,Q_D^-)
 \leq\Pr(F_n^c)=o(1)\leq\eta,
 \qquad \Delta_D=2c_DLw_n^2\asymp Ld^{-2}.
\tag{33}
\]

We reduce these pairs to the two minimax criteria.  For any pair with targets \(\theta^+,\theta^-\), laws \(Q^+,Q^-\), and separation \(\Delta\), define \(A_T=\{|T-\theta^+|\geq\Delta/2\}\).  On \(A_T^c\), the triangle inequality gives \(|T-\theta^-|\geq\Delta/2\), and hence
\[
\begin{aligned}
 \mathbb E_{Q^+}|T-\theta^+|+\mathbb E_{Q^-}|T-\theta^-|
 &\geq\frac\Delta2\{Q^+(A_T)+Q^-(A_T^c)\}\\
 &\geq\frac\Delta2\{1-\operatorname{TV}(Q^+,Q^-)\}.
\end{aligned}
\tag{34}
\]
Thus the larger absolute-error risk is at least \(\Delta(1-\eta)/4\).

For honest expected length, \({\rm def:minimax\text{-}criteria}\) permits data-measurable closed intervals \(\mathcal C(D)=[\ell(D),u(D)]\) with finite ordered endpoints.  Let such an interval satisfy the uniform coverage constraint.  Membership of \(Q^+\) and \(Q^-\) in \(\mathcal M_{n,d,L}\), coverage, and the definition of total variation give
\[
 Q^+(\theta^+\in\mathcal C)\geq1-\alpha,
 \qquad
 Q^+(\theta^-\in\mathcal C)
 \geq1-\alpha-\operatorname{TV}(Q^+,Q^-).
\]
The union bound therefore yields
\[
 Q^+\{\theta^+,\theta^-\in\mathcal C\}
 \geq1-2\alpha-\operatorname{TV}(Q^+,Q^-)
 \geq(1-2\alpha)/2.
\tag{35}
\]
On this simultaneous-inclusion event, connectedness and endpoint ordering force
\[
 \operatorname{length}\{\mathcal C(D)\}=u(D)-\ell(D)
 \geq|\theta^+-\theta^-|=\Delta.
\]
Therefore \(\mathbb E_{Q^+}\operatorname{length}(\mathcal C)\geq\Delta(1-2\alpha)/2\).  The interval restriction is used here explicitly; this implication is not asserted for arbitrary disconnected confidence sets.  Equations (25), (30), and (33)--(35), followed by the infima in \({\rm def:minimax\text{-}criteria}\), lower-bound both criteria by fixed multiples of the graph, interior, and lattice separations in their respective regimes.

It remains to compare those separations with the full supported envelope.  Put \(a_n=(n\sqrt d)^{-1/2}\).  For \(L>0\), (29) gives
\[
 Lh_0^2=a_nh_0^{-3/2}=L^{3/7}(n^2d)^{-1/7}.
\tag{36}
\]
If \(h_0\leq d^{-1}\), then \(L^2\geq d^{13/2}/n\), so \(a_nd^{3/2}\leq Ld^{-2}\); evaluating at \(h=d^{-1}\) gives
\[
 \inf_{h\in\mathcal H_{n,d}}
 \{Lh^2+a_nh^{-3/2}\}\leq2Ld^{-2}.
\tag{37}
\]
If \(d^{-1}<h_0<d^{-1/2}\), evaluation at \(h_0\) gives
\[
 \inf_{h\in\mathcal H_{n,d}}
 \{Lh^2+a_nh^{-3/2}\}
 \leq2L^{3/7}(n^2d)^{-1/7}.
\tag{38}
\]
If \(h_0\geq d^{-1/2}\), then \(L^2\leq d^3/n\), so \(Ld^{-1}\leq\sqrt{d/n}\); evaluation at \(h=d^{-1/2}\) gives
\[
 \inf_{h\in\mathcal H_{n,d}}
 \{Lh^2+a_nh^{-3/2}\}\leq2\sqrt{d/n}.
\tag{39}
\]
For \(L=0\), the infimum equals \(\sqrt{d/n}\).  Combine the graph and lattice pairs in (37), the graph and interior pairs in (38), and the graph pair alone in (39) and when \(L=0\).  A quantity bounded below by fixed multiples of nonnegative \(x\) and \(y\) is bounded below by a fixed multiple of \(x+y\).  By \({\rm def:frontier\text{-}envelope}\), both minimax criteria consequently satisfy
\[
 R^*_{n,d,L}\geq c\mathfrak r^{\mathrm{sup}}_{n,d,L},
 \qquad
 \Lambda^*_{n,d,L}(\alpha)
 \geq c\mathfrak r^{\mathrm{sup}}_{n,d,L}
\tag{40}
\]
for all sufficiently large \(n\).  Every response construction retained the arbitrary admissible kernel and the independent uniform latent types.  All divisions by \(B\), \(d\), or \(N_i\) are justified respectively by \(B>0\), \(d>0\), and restriction to \(N_i>0\); nonlinear constructions are used only when \(L>0\).  This proves the unrestricted converse on every eventually nonempty repaired class under the corrected nonempty-domain definitions.